% p3-07-mdl

\section{P3 §7. MDL Interpretation of Representation Quality}

7. MDL Interpretation of Representation Quality

  Figure (fig-p3-ch7-mdl). MDL-interpretation triptych --- Two-part code / Algorithmic information /
                                        42-entry catalogue.

  7.1 The Two-Part MDL Code
  The MDL principle evaluates a model by the total code length required to transmit both
  the model and the data given the model. In the GOLDEN BRIDGE (Trinity-Pellis Bridge) context:

  - Model: Trinity monomial T^ = PiL(x), requiring ell(lambda^) bits (Definition 2.6);
  - Data given model: Pellis residual c res of depth D, requiring approximately D $\cdot$ H(c res)
  bits, where H is the empirical entropy of the coefficient sequence.

  The two-part MDL framework --- in which model and residuals are coded separately and
  combined --- is the formulation of Barron, Rissanen, and Yu (1998). The full refined
  MDL treatment, including normalized maximum likelihood, is covered by Grünwald
  (2007, MIT Press).

        Definition 7.1 (MDL score).

  MDL(T^, c res) = ell(lambda^) + sumk=0D ell(ck res),

  where ell(c) = lceil log2(|c|+1) rceil + 1 for each coefficient (one bit for the sign).

        Theorem 7.2 (MDL optimality; Theorem). Among all representations of x as T
        + r with T $\in$ T and r approximated to precision $\varepsilon$ by a Pellis expansion of depth D,
        the representation (PiL(x), RL(x)) minimises the MDL score (Definition 7.1) up to
        an additive constant of O(log log(1/$\varepsilon$)).

  Proof sketch. Among all T $\in$ T with ell(lambdaT) $\leq$ L, the choice T^* = PiL(x) minimises
  |r| = |x - T|. A smaller residual |r| yields shorter Pellis coefficient sequences (since RD $\leq
  $\varphi$-(D+1) and the coefficients are bounded by the residual magnitude), hence a lower
  MDL score. The O(loglog(1/$\varepsilon$)) additive error arises from the integer rounding in
  coefficient coding. □

  7.2 Connection to Algorithmic Information Theory
  The relation between MDL scores and Kolmogorov/algorithmic complexity was
  established by Chaitin (1987). In the GOLDEN BRIDGE (Trinity-Pellis Bridge) setting:

  K(x|$\varepsilon$) $\leq$ MDL(T^*, c res) + O(log K(x|$\varepsilon$)),

  where K(x|$\varepsilon$) is the Kolmogorov complexity of x given precision $\varepsilon$. Constants with
  anomalously low MDL scores relative to their Kolmogorov complexity --- i.e., those that
  admit very short GOLDEN BRIDGE (Trinity-Pellis Bridge) descriptions compared to their informational content ---
  are candidates for genuine structural relationships with the Trinity monomial system.

  As Chaitin's framework makes clear, low algorithmic complexity is a necessary but not
  sufficient condition for structural understanding: it shows that a short description
  exists, not that the description reveals a physical law. This caveat is central to the
  falsification protocol of Section 9.4.

  7.3 MDL and the 42-Entry Catalogue
  The PhD programme's 42-entry precision catalogue provides an empirically
  pre-selected set of target constants together with their proposed symbolic
  representations. In the Paper 3 framework, each catalogue entry (xi, Fi) --- where Fi is
  the proposed symbolic form --- is reinterpreted as a specific Trinity label lambdai
  together with a Pellis residual ci res. The MDL score of this representation is then a
  computable compression metric.

        Definition 7.3 (Catalogue MDL vector). Let {(xi, lambdai)}i=142 be the 42
        catalogue entries. The catalogue MDL vector is

  M = (MDL(T(lambdai), ci res))i=142 $\in$ R42.

  The distribution of M over the catalogue --- its mean, variance, and outlier structure ---
  is the primary empirical object of the computational program (Section 9).

  Caveat. The 42 catalogue entries were assembled as an empirical survey of notable
  numerical coincidences and are explicitly labelled as not constituting proofs of
  structural relationships. The MDL reinterpretation preserves this status: it converts
  each entry from a proposed symbolic form to a compression score, but the score alone
  does not constitute a proof of any underlying mathematical or physical relationship.

  The distribution of MDL scores admits a natural three-way classification:

  Here L decimal is the description length of the naive decimal expansion at the same
  precision. High-MDL entries in the catalogue are candidates for removal or revision.

\subsection{References}

- Vasilev, D. \textit{Flos Aureus Monograph: Trinity Constants}. DOI \href{https://zenodo.org/doi/10.5281/zenodo.19227877}{10.5281/zenodo.19227877}.
- CODATA 2018 recommended values: NIST \href{https://physics.nist.gov/cuu/Constants/}{https://physics.nist.gov/cuu/Constants/}.
- Heyrovska, R. Golden-angle FSC publications --- Heyrovský Institute of Physical Chemistry.
